\chapter{Polysemy}
\section{Word Sense Induction}
\label{polysemyProblemForSemanticRelations}
In order to reach that point, we first need to define the basic object to be acquired, i.e. semantic relations. Unfortunately, this is not a precisely defined research object per se because it strongly depends on our previous knowledge about the possible word senses \parencite[12]{ayseExtractionSemanticWord2011}. Consequently, the sense of a word has to be defined before we can start to learn more abstract representations like hyponymy \parencite[137]{bartunovBreakingSticksAmbiguities2016}. Further obstacles may occur in the case of specific relations like synonymy, where a good operationalization is needed to retrieve the relevant instances from a corpus of text data \parencite[274]{divjakCorpusbasedCognitiveSemantics2009}. In my study, I will focus heavily on the approach of distributional semantics, describing specific semantic relations between words in terms of their co-occurrences, i.e. typical contexts in a corpus \parencites[162]{harrisDistributionalStructure1954}[30]{firthSynopsisLinguisticTheory1957}. This approach has traditionally suffered from several limitations. In some cases (e.g. synonymy), building unified representations for every type in a corpus can prove highly problematic \parencite[114]{karanDistributionalSemanticsApproach2012}. However, this is not universally true anymore as there are \gls{DSM}s that consistently distinguish between various contexts (\hyperlink{ELMO}{see below}). Other problems of distributional semantics are tied to fine-grained distinctions between specific semantic relations \parencite[115]{karanDistributionalSemanticsApproach2012}.

Co-occurrence counts can be a good basis for the detection of semantic relations in general, but they are rarely sufficient for a thorough classification into several subcategories. Nonetheless, previous research has been carried out (more or less successfully) using, e.g., association measures to detect synonymy \parencite[566]{hagiwaraSupervisedSynonymAcquisition2009}. One of the critical parts in that regard is the term `word pair', which introduces the linguistically ill-defined concept of `word' into an otherwise seemingly precise mathematical formula. In the following, we will start by using the term `word' simply as `token in a corpus' and later proceed to a more advanced definition. This progression seems reasonable given the fact that the usage of token-based \gls{PMI} for the research purposes of distributional semantics often leads to severely skewed results \parencite[444]{herbelotMeasuringSemanticContent2013}. Therefore, corpora may have to be tokenized and segmented using more complex separator rules than `whitespace and punctuation' before they can serve as a basis for distributional analyses \parencite[6]{mikolovDistributedRepresentationsWords2013}. Besides, polysemy may also become a problem for resolving other issues like synonymy, e.g. if two words are synonymous only in one of their possibly numerous senses \parencite[444]{herbelotMeasuringSemanticContent2013}. 

In addition to polysemy itself, related concepts like homonymy can also become obstacles with regard to the aim of defining semantic relations precisely \parencite[251]{gardnerValidatingConstructWord2007}. Distributional semantics provides a valuable framework to break down the multiple aspects of meaning for homonymous or polysemous words into smaller subsets, based on the combination with, e.g., adjectival modifiers \parencite[42]{boledaIntensionalityWasOnly2013}. Therefore, the integration of specific contextual elements into the analysis of hyponymy between Latin nouns seems to be a reasonable basis for distinguishing subtle semantic nuances more consistently. As a consequence, polysemy will be treated in this study not just as one of the analytical targets, but also as a methodological aspect that needs to be addressed before other kinds of analyses can be performed. A useful operationalization of polysemy in a distributional context is to look at the collocational diversity \parencite[8]{hamiltonDiachronicWordEmbeddings2016}. This approach brings along the problem of sampling, which is a general one for the whole of corpus linguistics: If we want to model a specific language like Latin as a whole but only use a tiny subset of the available language data for our corpus, how can we expect our results to be applicable to other Latin texts? The key point here is the purpose: For this study, our ultimate goal is not to model Latin nouns per se, but only with regard to the rather well-defined application context of language learning. Therefore, the explanatory power of our linguistic model does not need to be as comprehensive, and it will be restricted to predefined parameters. These parameters may in turn be described as, e.g., the objective of maximizing a student's learning success by providing structured information about semantic relations.

%TODO: add chapter on the choice of PROIEL as corpus... why not AGLDT or ITTB?

\section{Distributional Semantic Models for Word Senses}
Historically, linguistic features of words have often been represented in \gls{CVS}s \parencites{mikolovEfficientEstimationWord2013}{daiDocumentEmbeddingParagraph2015}{akbikContextualStringEmbeddings2018} that can be seen as statistical models of word distributions in natural language \parencite[1140]{bengioNeuralProbabilisticLanguage2003}. Such models map single elements of an utterance (e.g. characters, tokens, types, phrases etc.) to vectors in an n-dimensional \gls{CVS}. Depending on the dimensionality, various kinds of information will be integrated into the model. For instance, lexically related words (e.g. \textit{amor} `love' in PROIEL 86055.22\footnote{For details on the PROIEL corpus, see \textcite{haugCreatingParallelTreebank2008}. Note that, in this study, we only consider the Latin part of the corpus. Dot-separated unique identifiers will be used to refer to segments and tokens in the PROIEL corpus. This is often more precise than the conventional citation model for ancient works (in this case: Cic. off. 1.12) and does not require knowledge about abbreviations of author names or work titles. On the downside, it is corpus-specific and does not offer immediate information about author, work and structural position of the text passage.} and \textit{amare} `to love' in PROIEL 86334.11) usually correspond to vectors whose distance in the \gls{CVS} is rather small. Since most of these models try to predict words in a corpus given their surrounding context \parencite[5]{mikolovEfficientEstimationWord2013}, those words that co-occur frequently will also be mapped to similar vectors in the \gls{CVS}. This basic interaction helps us to distinguish multiple senses of a polysemous word empirically: If each word sense is associated with different contexts and, thus, different co-occurring words, its calculated vector will be different from those of other word senses. Therefore, a very simple evaluation for \gls{DSM}s with regard to \gls{WSI} looks like this:
\begin{itemize}
	\item Given a pair of contexts that
	\item both contain the same word,
	\item classify whether that target word is used in the same sense or not.
\end{itemize}
Formalizing this into a dataset structure, we may provide a plain text file where each each line contains one example for classification. Since Latin is a highly inflectional language, morphological features of the targeted word can differ a lot. Therefore, the actual word forms have to be indicated separately for each context in the example. If they occur multiple times in a context, the dataset should specify which occurrence of the word is to be evaluated. Finally, a binary value (such as `0 / 1' or `true / false') indicates the ground truth for this classification example, i.e. whether the word senses are the same or not. This is the most critical part: `Ground truth' is a dangerous assumption for linguistic phenomena because it suggests a consistent, unchanging goal to be achieved \parencite[129]{leeGeneratingGroundTruth2012}. This desire stems from the formalization that is deemed necessary to systematically evaluate computational language models \parencite[4]{fischerGroundTruthCreation2010}. Usually, such gold standards depend on human expertise, i.e. people that are highly acquainted with the target domain \parencite[519]{uzunerCommunityAnnotationExperiment2010}. In the case of historical languages like Latin, where specialized software and datasets are scarce \parencites[168]{karakantaNeuralMachineTranslation2018}[3]{mcgillivrayMethodsLatinComputational2013}, researchers often have to resort to digitizing analog resources or creating entirely new ones. This kind of curation process can and should rely on successful patterns from datasets for modern languages. Thus, it is reasonable to have a short look at existing comparable resources for evaluation.

The SemCor \parencite{mihalceaSemCorCorpus2008} and MASC-News \parencite[4218]{moroAnnotatingMASCCorpus2014} datasets are \gls{POS}-tagged English corpora where word senses have been annotated using WordNet \parencite{tengiDesignImplementationWordNet1998}. In that sense, such data is both empirical (because it contains authentic historical utterances) and intuition-based (because it uses an external knowledge base for defining the word senses, \hyperref[WordNet]{see below}). There have been attempts to unify \parencite[216]{navigliBabelNetBuildingVery2010} the various knowledge bases that researchers employed to annotate word senses. Promising results were achieved by discussing discrepancies in the face of text corpora, i.e. actual text passages from authentic sources \parencite[127]{bakerWordNetFrameNetComplementary2009}. This development highlights the importance of usage-based linguistics and points towards the need for a more rigorous, bottom-up formalization of word senses. Further approaches to creating datasets for \gls{WSD} include the usage of dictionaries (other than WordNet) as lexical resources \parencites[25]{mihalceaSenseval3EnglishLexical2004}[2]{jimeno-yepesExploitingMeSHIndexing2011}[6]{yuanSemisupervisedWordSense2016} and the usage of glosses as representations of word senses \parencite[8]{luoIncorporatingGlossesNeural2018}. In contrast, \gls{WSI} usually does not rely on external sense inventories, but tries to infer word senses in an unsupervised, empirical manner only from the text itself. State-of-the-art models cluster distributions of vector representations for possible context-dependent substitutes of a word and, thus, are able to detect distinctions between word senses that had not previously been considered in intuition-based sense inventories \parencite[1-4]{amramiBetterSubstitutionbasedWord2019}. Predecessors to this approach directly used a word's surrounding context to construct sense representations for clustering and identified substantial quality criteria for the evaluation of \gls{WSI} models \parencite[876-878]{huangImprovingWordRepresentations2012}:
\begin{itemize}
	\item human judgment on word sense similarity
	\item variation in patterns of meaning (homonymy, polysemy etc.)
	\item variation in parts of speech (esp. nouns, verbs, adjectives)
\end{itemize}
We can certainly add that, in addition to variation, balancing is a decisive factor for the training data \parencites[620]{liuCombiningIntegratedSampling2011}[9]{weiRoleBalancedTraining2013}[1127]{boneApplyingMachineLearning2015}: If we are interested in rare linguistic phenomena, such as subtle semantic differences between almost synonymous words, we need to make sure that they are represented frequently in our training data. The third criterion is most controversial: On the one hand, powerful \gls{WSI} models need to be general enough to provide sufficient results under different circumstances, e.g. in varying domains and registers. On the other hand, modeling word senses for different parts of speech needs to take into account the diverse influences of lexis and syntax: In some settings, nouns may be represented as sets of their hyponyms \parencite[640]{lewisCompositionalHyponymyPositive2019}, while transitive verbs can be modeled as triples together with their arguments, e.g. subjects and objects \parencite[9]{grefenstetteExperimentalSupportCategorical2011}. Such different approaches uncover the intricate nature of formalizing fine-grained semantic distinctions from a distributional perspective. As a consequence, our goal of detecting hyponyms for language learning seems to warrant a modular operationalization with regard to parts of speech. 